\subsection{Exercices}

% --- Section 1 : Définitions, Skewness et Kurtosis ---

\begin{exercicebox}[Exercice 1 : Concepts (Moments)]
\begin{enumerate}
    \item À quoi correspond le 1er moment non centré, $E[X^1]$ ?
    \item À quoi correspond le 2ème moment centré, $E[(X-\mu)^2]$ ?
\end{enumerate}
\end{exercicebox}

\begin{exercicebox}[Exercice 2 : Interprétation (Skewness)]
Une distribution des salaires dans une entreprise a un skewness de +2.5. Qu'est-ce que cela signifie concrètement sur la répartition des salaires ?
\end{exercicebox}

\begin{exercicebox}[Exercice 3 : Interprétation (Skewness)]
Si une distribution est parfaitement symétrique, que vaut son skewness ?
\end{exercicebox}

\begin{exercicebox}[Exercice 4 : Interprétation (Kurtosis)]
Une distribution des rendements d'un actif financier a un "excess kurtosis" de +5.0.
\begin{enumerate}
    \item Comment appelle-t-on ce type de distribution ?
    \item Qu'est-ce que cela implique sur la probabilité des "krachs" (événements extrêmes) par rapport à une loi normale ?
\end{enumerate}
\end{exercicebox}

\begin{exercicebox}[Exercice 5 : Interprétation (Kurtosis)]
Une distribution a un "excess kurtosis" de -1.0.
\begin{enumerate}
    \item Comment appelle-t-on ce type de distribution ?
    \item Comment décririez-vous visuellement son "pic" et ses "queues" par rapport à une loi normale ?
\end{enumerate}
\end{exercicebox}

\begin{exercicebox}[Exercice 6 : Comparaison (Skewness)]
La distribution exponentielle est-elle symétrique, asymétrique à gauche ou asymétrique à droite ? Quel est le signe de son skewness ?
\end{exercicebox}

\begin{exercicebox}[Exercice 7 : Comparaison (Kurtosis)]
La distribution uniforme est-elle mésokurtique, leptokurtique ou platykurtique ? Son "excess kurtosis" est-il positif, négatif ou nul ?
\end{exercicebox}

% --- Section 2 : Symétrie ---

\begin{exercicebox}[Exercice 8 : Propriété de Symétrie]
Soit $X$ une variable aléatoire symétrique autour de sa moyenne $\mu$.
Montrez que son 3ème moment centré, $E[(X-\mu)^3]$, est nul (et donc que son skewness est nul).
\end{exercicebox}

\begin{exercicebox}[Exercice 9 : Vérification de Symétrie (PDF)]
La PDF d'une loi $X$ est $f(x) = \frac{1}{2}e^{-|x-3|}$.
Cette distribution est-elle symétrique ? Si oui, autour de quel point $\mu$ ?
\end{exercicebox}

\begin{exercicebox}[Exercice 10 : Symétrie et Moments]
Si une distribution a un skewness de 0, est-elle forcément symétrique ? (Indice : Pensez à une distribution qui aurait des queues asymétriques mais qui s'annuleraient pour le moment d'ordre 3).
\end{exercicebox}

% --- Section 3 : Moments d'Échantillon vs Population ---

\begin{exercicebox}[Exercice 11 : Définitions (Échantillon vs Pop.)]
Quelle est la différence conceptuelle entre $\sigma^2$ et $s^2$ ?
\end{exercicebox}

\begin{exercicebox}[Exercice 12 : Correction de Bessel ($n-1$)]
Pourquoi utilise-t-on $n-1$ au dénominateur pour $s^2$ ? Que se passerait-il si nous avions un échantillon de $n=1$ et que nous divisions par $n$ ?
\end{exercicebox}

\begin{exercicebox}[Exercice 13 : Calcul (Moments d'Échantillon)]
On observe l'échantillon de données suivant : $\{ 2, 3, 10 \}$.
\begin{enumerate}
    \item Calculez la moyenne d'échantillon $\bar{X}$.
    \item Calculez la variance d'échantillon non biaisée $s^2$.
\end{enumerate}
\end{exercicebox}

% --- Section 4 : Fonctions Génératrices (MGF) - Calculs ---

\begin{exercicebox}[Exercice 14 : MGF (Propriété de base)]
Soit $M_X(t)$ la MGF d'une variable $X$. Que vaut $M_X(0)$ ?
\end{exercicebox}

\begin{exercicebox}[Exercice 15 : MGF (Série de Taylor)]
Le développement en série de Taylor d'une MGF est $M_X(t) = 1 + 5t + 14t^2 + O(t^3)$.
(Rappel : $M_X(t) = 1 + E[X]t + E[X^2]\frac{t^2}{2!} + \dots$)
\begin{enumerate}
    \item Trouvez $E[X]$.
    \item Trouvez $E[X^2]$.
    \item Calculez $\text{Var}(X)$.
\end{enumerate}
\end{exercicebox}

\begin{exercicebox}[Exercice 16 : MGF (Calcul de moments)]
La MGF de la loi Exponentielle de paramètre $\lambda$ est $M_X(t) = \frac{\lambda}{\lambda - t}$.
\begin{enumerate}
    \item Calculez $M_X'(t)$ et trouvez $E[X] = M_X'(0)$.
    \item Calculez $M_X''(t)$ et trouvez $E[X^2] = M_X''(0)$.
\end{enumerate}
\end{exercicebox}

\begin{exercicebox}[Exercice 17 : MGF (Variance de l'Exponentielle)]
En utilisant les résultats de l'exercice 16, déduisez $\text{Var}(X)$ pour une loi Exponentielle($\lambda$).
\end{exercicebox}

\begin{exercicebox}[Exercice 18 : MGF (Loi Normale)]
La MGF de $X \sim \mathcal{N}(\mu, \sigma^2)$ est $M_X(t) = \exp(\mu t + \frac{1}{2}\sigma^2 t^2)$.
Calculez $M_X'(t)$ et vérifiez que $M_X'(0) = \mu$.
\end{exercicebox}

\begin{exercicebox}[Exercice 19 : MGF (Loi Normale)]
En utilisant la MGF de l'exercice 18, calculez $M_X''(t)$ et montrez que $E[X^2] = \mu^2 + \sigma^2$.
\end{exercicebox}

\begin{exercicebox}[Exercice 20 : MGF (Loi Normale)]
En utilisant les résultats des exercices 18 et 19, retrouvez la formule de la variance $\text{Var}(X)$.
\end{exercicebox}

% --- Section 5 : MGF - Sommes de Variables Indépendantes ---

\begin{exercicebox}[Exercice 21 : Propriété des Sommes]
Soient $X$ et $Y$ deux variables aléatoires \textbf{indépendantes}. Soit $S = X+Y$.
Comment la MGF de $S$, $M_S(t)$, est-elle liée à $M_X(t)$ et $M_Y(t)$ ?
\end{exercicebox}

\begin{exercicebox}[Exercice 22 : Application (Somme de Poissons)]
$X \sim \text{Poisson}(\lambda_1)$ a pour MGF $M_X(t) = e^{\lambda_1(e^t - 1)}$.
$Y \sim \text{Poisson}(\lambda_2)$ a pour MGF $M_Y(t) = e^{\lambda_2(e^t - 1)}$.
$X$ et $Y$ sont indépendantes.
\begin{enumerate}
    \item Trouvez la MGF de $S = X+Y$.
    \item En reconnaissant la forme de la MGF de $S$, quelle est la loi de $S$ ?
\end{enumerate}
\end{exercicebox}

\begin{exercicebox}[Exercice 23 : Application (Somme de Binomiales)]
La MGF de $X \sim \text{Bin}(n, p)$ est $M_X(t) = (p e^t + (1-p))^n$.
Soit $Y \sim \text{Bin}(m, p)$ (même $p$), indépendante de $X$.
Quelle est la loi de $S = X+Y$ ? Justifiez avec les MGF.
\end{exercicebox}

\begin{exercicebox}[Exercice 24 : Application (Transformation Linéaire)]
Soit $M_X(t)$ la MGF de $X$. Soit $Y = aX + b$.
Exprimez $M_Y(t)$ en fonction de $M_X(t)$.
(Indice : $M_Y(t) = E[e^{t(aX+b)}]$).
\end{exercicebox}

\begin{exercicebox}[Exercice 25 : Application (Moyenne d'Échantillon)]
Soient $X_1, \dots, X_n$ des v.a. i.i.d. (indépendantes et identiquement distribuées) avec la MGF $M_X(t)$.
Soit $\bar{X} = \frac{1}{n} \sum X_i$ la moyenne d'échantillon.
Exprimez la MGF de $\bar{X}$ en fonction de $M_X(t)$.
(Indice : utilisez les résultats des exercices 21 et 24).
\end{exercicebox}